\chapter{NixOS 的理由}
\label{ch:why-nixos}
\chapterlead{本章说明系统状态从哪里来，这个来源如何决定复现、回退与审计的成本。Windows、macOS 与 Linux 传统发行版在这一点上属于同一类，NixOS 属于另一类。后半部分把这些性质放进 AI Agent 参与系统维护的场景里评估，并给出可核查的规模数字与明确的边界。}

\section{起步：清点一台机器的现状}

在纸上回答四个问题。

\begin{enumerate}
  \item 这台机器装了哪些软件？
  \item 每个软件的版本是多少？
  \item 哪些配置文件被改过，改成了什么？
  \item 换一台空机器，恢复到现在的状态需要多久？
\end{enumerate}

前两个问题有部分答案。Windows 的“应用和功能”、macOS 的“应用程序”文件夹、Linux 的包管理器都能列出软件名，版本要在每个程序的“关于”里逐个查看。第三个问题通常没有答案：配置散落在注册表、plist、\texttt{/etc}、用户主目录与各程序自己的目录里，改过哪些依赖记忆。

四个问题合起来是系统的状态。第四个问题的答案决定一次故障的代价：重装一遍要一天，还是一小时。

\section{累积的系统}

读者手上的机器多半是 Windows 或 macOS。这两种系统的状态来源，与 Linux 传统发行版相同。

Windows 的安装与配置是一串操作。双击安装包，程序把文件写进 \texttt{Program Files}，把设置写进注册表，把用户数据写进 \texttt{AppData}。卸载时删除的多是安装器登记过的那部分；注册表里的残留、\texttt{AppData} 下的缓存、别的程序建立的关联项都留在原处。安装顺序也留下痕迹：先装 A 再装 B，与先装 B 再装 A，结果可能不同。

macOS 的目录结构整齐，来源相同。应用是 \texttt{/Applications} 下的包，系统级设置散在 \texttt{/Library} 与 \texttt{/etc}，用户级设置散在 \texttt{\textasciitilde/Library} 下的 plist 文件里，开机启动项是 \texttt{launchd} 的配置文件。安装开发工具常要再执行几条命令，把文件放进 \texttt{/usr/local} 或 \texttt{/opt/homebrew}。半年之后，某个 plist 由谁写入已无从查起。

Linux 传统发行版的形态还不熟悉，机制与上面两者相同：执行 \texttt{apt install} 装一个包，编辑 \texttt{/etc} 下的文件改一项设置，添加一个第三方软件源，半年做一次版本升级。

四种系统的对照如下表。前三行都是命令式系统，状态的最终形态由操作历史决定。

\ifbknarrow
\begin{itemize}
  \item \textbf{Windows}：状态由注册表、\texttt{Program Files} 与 \texttt{AppData} 逐次累积；重装后按记忆重配，\texttt{winget} 只列出包名；系统还原退回某一时刻的系统文件与注册表，个人文件不在其中。
  \item \textbf{macOS}：状态由 \texttt{/Library} 与 \texttt{\textasciitilde/Library} 的 plist、\texttt{launchd} 配置与应用包累积；迁移助理复制旧机器，Brewfile 只列出包名；Time Machine 以备份为单位，前提是备份已经在手。
  \item \textbf{Linux 传统发行版}：状态由 \texttt{/etc}、包管理器与 dotfiles 累积；复现靠照抄步骤；文件系统快照回到某一时刻的磁盘。
  \item \textbf{NixOS}：状态由配置表达式与锁文件算出；复现靠同一组输入重建；代回到上一次构建的结果。
\end{itemize}
\else
\begin{center}
\centering
\small
\begin{booktable}{\textwidth}{L{0.80} L{1.20} L{0.95} J{1.05}}
\caption{四种系统的状态来源与回退方式}\\
\toprule
系统 & 状态来源 & 复现方式 & 回退方式 \\
\midrule
\endfirsthead
\toprule
系统 & 状态来源 & 复现方式 & 回退方式 \\
\midrule
\endhead
Windows & 注册表、\texttt{Program Files} 与 \texttt{AppData} 逐次累积 & 重装后按记忆重配；\texttt{winget} 只列出包名 & 系统还原：某一时刻的系统文件与注册表 \\
macOS & \texttt{/Library} 与 \texttt{\textasciitilde/Library} 的 plist、\texttt{launchd} 配置、应用包 & 迁移助理复制旧机；Brewfile 只列出包名 & Time Machine：备份，须事先存在 \\
Linux 传统发行版 & \texttt{/etc}、包管理器、dotfiles 逐次累积 & 照抄安装步骤 & 文件系统快照：某一时刻的磁盘 \\
NixOS & 配置表达式与锁文件 & 同一组输入重建 & 代：上一次构建的结果 \\
\bottomrule
\end{booktable}
\end{center}
\fi

三条后果在 Windows、macOS 与 Linux 传统发行版上都能观察到。

第一，状态没有单一出处。系统的当前形态由操作历史决定，历史分散在注册表、plist、\texttt{/etc}、命令记录与人的记忆里。同一份安装文档在两台机器上执行，得到的往往不是同一套软件：源里的版本已经变了，某一步依赖前一步的残留。

第二，复现只能照抄步骤。把步骤写成脚本可以减轻问题，脚本本身又需要维护；脚本执行失败时，系统停在中途状态，无法判断哪些步骤已经生效。Windows 的 \texttt{winget} 与 macOS 的 Homebrew 能把包名列成一份清单，清单不带版本与构建输入，换一台机器装出来的是当天的最新版。

第三，回退需要事先准备。Windows 的系统还原、macOS 的 Time Machine、Linux 的文件系统快照都只能退回某个时刻的磁盘内容，前提是那个时刻的记录已经存在。三者回答“升级前是什么样”，不回答“这套系统由哪些输入构建”。

这类系统称为命令式（imperative）：使用者给出改变状态的命令，系统执行命令，状态随之改变。界面与包格式各不相同，状态的来源是同一个。

\begin{culturenote}
声明式配置管理并不只属于 NixOS。Windows 的组策略与 Intune、macOS 的 MDM、Linux 上的 CFEngine、Puppet、Ansible 同样用配置文件描述目标状态，差别在作用范围。它们把一台已经存在的机器收敛到描述的状态，机器上仍保留未被描述的部分；NixOS 用描述直接构建出系统，未被描述的部分不属于系统。前者的结果是“看起来一样”，后者的结果是“由同一组输入算出”。
\end{culturenote}

\section{生成的系统}

NixOS 把系统写成表达式。配置文件声明装哪些包、开哪些服务、用户与网络如何设置。构建系统读取这份表达式，算出一个完整的系统，切换后生效。

一次重建分四个阶段，每个阶段都有可观察的产物。

\begin{itemize}
  \item 求值：模块系统合并配置，检查选项名与类型。写错选项名在这一步报错，系统尚未改动。
  \item 构建：按输入在沙箱里构建软件，产物落到 \texttt{/nix/store} 下以哈希命名的路径。
  \item 生成：把闭包与启动脚本组织成一个代（generation），写入引导菜单。
  \item 切换：激活新代。旧代保留在原处，随时可以选回。
\end{itemize}

flake 的锁文件固定外部输入的提交号。同一份配置与同一份锁文件，在同一体系结构的不同机器上算出同一组 store 路径。这是可复现的具体含义：判断依据是输入相同，结果相似不足以说明。

配置可以进版本控制。谁在什么时候改了哪一项，\texttt{git log} 有记录；改动是否成立，构建与测试给出答案。

这套写法不限于 NixOS。Nix 本身运行在 Linux 与 macOS 上，Windows 需经 WSL。nix-darwin 用同一套模块写法管理 macOS 的包、服务与系统默认值，本书不涉及。

\section{可复现的范围}

可复现指同一组输入得到同一组产物，不指整台机器逐字节相同。

由配置决定的部分可以复现：软件及其版本、服务配置、用户与权限、防火墙规则、桌面环境。这些内容写在表达式里，构建结果由哈希固定。

不由配置决定的部分不能复现：数据库里的数据、用户主目录、日志、密钥、模型权重、已经执行过的数据库迁移。重建系统不会恢复这些内容，回滚系统也不会撤销迁移。

还有一层更细的差别。同样的输入在不同机器上构建，产物可能不完全相同：时间戳、构建路径、并行顺序都会留下痕迹。NixOS 的可复现构建项目持续统计这类差异，目标是把它们逐个消除。要求逐字节相同时，需要在目标环境独立重建并比较，不能只凭用过 Nix 作判断。

\begin{notebox}[模型权重也要固定]
下载模型权重时按最新版的 URL 取文件，内容随时间变化，系统重建后行为随之改变。把权重的 SHA-256 写进配置，哈希不符时构建失败，模型版本便与系统代际一起回滚。本地推理部署常用这种做法，本书不展开。
\end{notebox}

\section{回滚的两种语义}

快照回滚撤销一次操作。系统在升级前拍一张磁盘快照，出问题时从引导菜单选回快照，磁盘回到拍摄时刻。这种做法适合撤销一次坏更新，范围是整个根文件系统，通常不包括用户数据。

代回滚回到一个构建结果。旧代是上一次构建的完整系统，选择它等于把系统切换回那组输入算出的状态。配置文件、包版本、服务定义一起回退，store 里的旧路径仍在原处，不需要重新下载。

两者的差别在可解释性。快照能回答“什么时候”，代能回答“由什么构建”。系统起不来时，两者都能救急；事后追究这次改动为何出错，只有后者留下完整的输入记录。

代价是磁盘。每个代引用的闭包都留在 store 里，路径在代际之间共享，但旧代本身不会自动删除。定期清理旧代是日常运维的一部分，命令见第\ref{ch:upgrade-rollback}章。

Windows 与 macOS 上更常见的两种回退也属于快照一类，范围值得辨清。系统还原保存某一时刻的系统文件、注册表与驱动，个人文件不在其中，功能在部分机器上默认关闭。Time Machine 以备份为单位，整盘恢复会把个人文件一并带回，备份之后新写的文件随之消失。两者都要求那个时刻的记录事先存在。

\section{代价}

收益需要与代价一起看。下列六项是使用 NixOS 的主要成本，括号内是 2026 年 11 月的量级。

\begin{enumerate}
  \item 学习成本。需要掌握一小部分 Nix 语言与模块规则。2025 年社区调查的 3,399 份完整回答中，28.9\% 使用不足一年，82.5\% 不足四年，多数使用者仍在爬坡。
  \item 文档。官方在 2026 年 7 月 10 日的公告中承认，新用户经常在文档里受挫，正在为文档团队募资。
  \item 求值时间。改动底层输入会触发大量重建。Nix 2.35 优化了输入抓取与求值性能，这项成本仍然存在。
  \item 磁盘。同一软件的多个版本并存，旧代保留。定期执行垃圾回收是必需的维护动作。
  \item flakes 未稳定。flake 自 2021 年作为实验特性引入，2026 年 10 月仍没有转正时间表，接口可能变化。
  \item 支持窗口与平台。26.05 的安全支持到 2026 年 12 月 31 日，约七个月；26.05 是最后一个支持 Intel Mac（x86\_64-darwin）的版本，26.11 起不再构建。
\end{enumerate}

从 Windows 或 macOS 迁入时，还有一项成本不在表内：终端与包管理的操作习惯需要重建。本书第一部分的顺序按这件事安排，第\ref{ch:first-terminal}章到第\ref{ch:network}章依次建立这套习惯。

\begin{warnbox}[未收录的软件要自行打包]
nixpkgs 收录十四万个包，仍有软件不在其中。自行编写打包表达式需要读上游的构建脚本，这是初学者最容易停住的一步。先用 \texttt{nix search} 确认有没有现成的包，再决定是否自己写。
\end{warnbox}

\section{AI 时代的需求变化}

Stack Overflow 2026 年的开发者调查给出五个数字：83\% 的受访者使用至少一种 AI 工具，65.9\% 使用编码助手，26.2\% 使用 agent 或自动化工作流，48\% 只在容易验证时才信任 AI 的输出，76.5\% 会先在本地运行 AI 给出的答案。

最后一个数字说明问题所在。生成的代价在下降，验证的代价没有随之下降。改动由 agent 提出时，提出十个方案很容易，判断哪一个满足约束才是瓶颈。

这改变了系统配置中三项性质的相对价值。

\begin{enumerate}
  \item 改动在生效前可被检查。选项名、类型、依赖、服务能否启动，都能在切换之前得到答案。
  \item 失败的改动可以整次退回。退回的对象是一份构建结果，不是一次操作。
  \item 环境可以作为输入交付。配置与锁文件交给另一台机器或另一个 agent，得到同一组软件。
\end{enumerate}

三项性质在传统发行版上也能部分实现，代价是额外的工具与约定；NixOS 把这三项放进系统本身。

\begin{figure}[H]
\centering
\figfit{%
\begin{tikzpicture}[font=\small, node distance=4mm,
  gate/.style={draw=ink!45, fill=ink!5, rounded corners=1pt, inner sep=3pt, align=center, minimum width=74mm, minimum height=9mm}]
  \node[gate] (a) {求值\\[-1pt]{\scriptsize\color{muted}选项名与类型检查}};
  \node[gate, below=of a] (b) {构建\\[-1pt]{\scriptsize\color{muted}沙箱内按输入算出 store 路径}};
  \node[gate, below=of b] (c) {虚拟机测试\\[-1pt]{\scriptsize\color{muted}启动一台机器，检查服务是否应答}};
  \node[gate, below=of c] (d) {切换\\[-1pt]{\scriptsize\color{muted}新代写入引导菜单}};
  \node[gate, below=of d, draw=teal!55, fill=teal!8] (e) {回退\\[-1pt]{\scriptsize\color{muted}选择上一代，系统回到上一次构建}};
  \draw[-{Latex[length=1.5mm]}, teal, thick] (a) -- (b);
  \draw[-{Latex[length=1.5mm]}, teal, thick] (b) -- (c);
  \draw[-{Latex[length=1.5mm]}, teal, thick] (c) -- (d);
  \draw[-{Latex[length=1.5mm]}, teal, thick, dashed] (d) -- (e);
\end{tikzpicture}%
}
\caption{一次系统变更经过的三道检查、一次切换与一条退路}
\end{figure}

\section{Agent 与系统配置的接触面}

系统本身不提供 agent 界面。要改系统，agent 仍然读配置、跑 \texttt{nixos-rebuild}，通常需要 \texttt{sudo}。它提供的是另一种接口：变更写成文本，构建给出成败，代给出退路。

三道检查对 agent 的作用各不相同。选项名写错，求值阶段报错；依赖缺失或构建脚本有问题，构建失败；服务配置错误但构建通过，虚拟机测试能暴露一部分。三者都通过而运行期仍然出错时，回退只需重启并选上一代。

读取方式也与查手册不同。选项叫什么、默认值是多少，答案在 nixpkgs 的模块源码里，比手册更全、更新更快。这项能力有代价：定位、版本对齐与上下文管理都是工程，读错版本、读漏调用点都会发生，静默失败仍要靠实际运行与人工验收。

\subsection{两项弱证据}

一项公开的基准测试记录了语料设计。3.0 版的 30 个任务里有 24 个被前沿模型凭记忆解出；3.1 版里凡求值阶段报错的缺陷都能解；唯一全部失败的任务是运行期静默失败，服务起不来而配置看不出问题。该基准是个人项目，仓库只有 3 个 star，两版语料的要求也不同，只能作为现象记录，不能当作通过率证据。

一份单人维护的机群记录（2026 年 11 月读数）显示：八台 NixOS 主机与一台 macOS 由同一个 flake 管理，共 151 个 \texttt{.nix} 文件、11,681 行。维护者自述不熟悉 Nix 语言，工作方式是提出要求并审查构建结果；仓库里另有两千余行面向 agent 的文档与 25 份决策记录。这是单一操作者的观察，不能外推为普遍情况。

由此得到的结论是条件句。有能干的 agent、有可以构建与回滚的真实机器、有人做最终判断时，NixOS 的上手门槛低于文档给人的印象；缺少这三项条件时，门槛不变。

边界还有三条。NixOS 处在构建与部署层，训练调度由 Kubernetes 与 Slurm 一类系统负责，它们在 NixOS 上作为服务运行。GPU 软件栈由社区维护，CUDA 团队 2022 年成立，五名成员并在公开寻求赞助；需要特定驱动版本或专有组件时仍要自行打包。只固定单个项目的工具链时，容器与 devcontainer 的学习与迁移成本通常更低，NixOS 的差别优势在整台机器与多台机器之间的一致性。

\section{规模与位置}

下列数字可以核查，日期与口径一并列出。

\begin{booktable}{\textwidth}{L{0.95} J{0.75} J{1.30}}
\toprule
指标 & 数值 & 口径与日期 \\
\midrule
\endhead
nixpkgs 包数 & 14 万以上 & nixos.org，2026-11 \\
26.05 贡献者 & 2,842 人 & 发布公告，2026-05-30 \\
26.05 提交数 & 59,703 条 & 同上，自 25.11 起 \\
二进制缓存 & 636 TB，10 亿对象 & 官方博客，2026-09-20 \\
社区调查 & 3,399 份完整回答 & 调查期 2025-09-01 至 12-15 \\
生产与 CI 使用 & 21.4\% 与 20.5\% & 同上，受访者自报 \\
DistroWatch & 第 20 位 & 页面点击，2026-11-16 \\
Stack Overflow & 587 个问题 & 累计，约为 archlinux 的 35\% \\
\bottomrule
\end{booktable}

这些数字说明三件事。第一，生态在扩张：贡献者数从 22.11 的 1,652 增长到 26.05 的 2,842，调查完整回答数从 2023 年的 1,708 增至 2025 年的 3,399；贡献者在 25.05 达到 2,857 之后进入平台期，26.05 为 2,842，包数量只有当期读数。第二，绝对规模仍然小：DistroWatch 的页面点击排在第 20 位，问答量约为 Ubuntu 的百分之一。第三，专业使用存在，但不是多数：受访者中两成把 Nix 用在生产服务器或 CI。

官方没有公布安装量或活跃用户数。上表的每一项都是代理指标，不能换算成有多少台机器装了 NixOS。机构采用的早期案例是荷兰政府的 DAWO 项目，以 NixOS 为基础构建公务工作环境，代码库到 2026 年 11 月仍有提交。关于试点规模，两份材料说法不一：一份写“小规模试点”，另一份称覆盖八个市镇；两者都来自二手报道，未见官方名单，全国部署时间表亦无一手材料。

\section{适用与不适用}

下表按使用场景给出判断，依据列说明判断的来源。

\begin{booktable}{\textwidth}{L{1.05} L{0.95} J{1.00}}
\toprule
情形 & 判断 & 依据 \\
\midrule
\endhead
学习 Linux、建立可回滚的实验环境 & 适用 & 试错成本低，配置可留存 \\
多台机器、需要一致配置 & 适用 & 一套模块覆盖全部主机 \\
需要逐项审计系统变更 & 适用 & 配置、构建、代三层记录 \\
机器只用一次、装完即弃 & 不适用 & 学习成本收不回来 \\
只固定单个项目的工具链 & 有条件 & 容器或 devcontainer 成本更低 \\
团队没有维护者与 agent & 不适用 & 文档门槛仍在 \\
需要厂商长期支持与合规背书 & 有条件 & 支持窗口约七个月，需商业支持 \\
日常办公、依赖 Windows 或 macOS 专有软件 & 不适用 & 专有软件不能脱离原系统 \\
已有 Mac、希望沿用同一套写法 & 有条件 & nix-darwin 可复用模块，不管理内核与硬件 \\
\bottomrule
\end{booktable}

本书的选择落在此表第一行。目标读者要建立的是一套可以反复改坏再修好的实验环境，回滚能力对这件事的作用大于上手速度。

\section{练习}

\begin{enumerate}
  \item 按本章开头的四个问题清点当前使用的机器，估计重装到可用状态需要的时间。
  \item 找出三项只在记忆里、没有写进任何配置文件的系统改动，写下它们的来源。
  \item 用一个段落说明快照回滚与代回滚的差别，各给一个适用场景。
  \item 用“代价”与“规模与位置”两节的数据，说明 NixOS 在增长、同时仍是小众发行版。
  \item 对比系统还原、Time Machine 与代回滚：各能回答哪些问题，各不能回答哪些问题。
  \item 列出一项本章没有提到、但 agent 修改系统时会需要的约束，说明理由。
\end{enumerate}

\begin{checkpoint}
本章结束时，应能：
\begin{itemize}
  \item 说明命令式与声明式两种系统状态的来源与后果；
  \item 用注册表、plist 与 \texttt{/etc} 各举一例，说明三种系统的状态来源相同；
  \item 描述一次重建的四个阶段，指出每个阶段的产物；
  \item 区分可复现覆盖的范围与不覆盖的范围；
  \item 说明快照回滚与代回滚各自的语义与代价；
  \item 列出使用 NixOS 的六项主要代价；
  \item 说明 AI 工具普及后，验证与回退为何比生成更重要；
  \item 用可核查的指标描述 NixOS 当前的规模与位置。
\end{itemize}
\end{checkpoint}
